\chapter{Chatbot Patterns}
\label{ch:chatbot}

One of the most common use cases for \haira{} is building chatbot backends. This chapter shows patterns from simple to advanced.

\section{Simple REST Chatbot}

A chatbot accessible via HTTP POST:

\begin{lstlisting}
import "io"
import "http"

provider anthropic {
    api_key: env("ANTHROPIC_API_KEY")
    model: "claude-sonnet-4-20250514"
}

agent Assistant {
    provider: anthropic
    system: "You are a helpful assistant."
    memory: conversation(max_turns: 50)
    temperature: 0.7
}

@webhook("/chat")
workflow Chat(message: string, session_id: string) -> { reply: string } {
    reply, err = Assistant.ask(message, session: session_id)
    if err != nil {
        return { reply: "Sorry, something went wrong." }
    }
    return { reply: reply }
}

fn main() {
    server = http.Server([Chat])
    io.println("Chat server on :8080")
    server.listen(8080)
}
\end{lstlisting}

Usage:

\begin{lstlisting}[language=bash,numbers=none]
$ curl -X POST http://localhost:8080/chat \
    -d '{"message": "Hello!", "session_id": "user-123"}'

{"reply": "Hi there! How can I help you today?"}

$ curl -X POST http://localhost:8080/chat \
    -d '{"message": "What did I just say?", "session_id": "user-123"}'

{"reply": "You said \"Hello!\""}
\end{lstlisting}

The agent remembers the conversation because both calls share the same \code{session\_id}.

\section{Streaming WebSocket Chatbot}

For real-time token-by-token responses over WebSocket:

\begin{lstlisting}
@websocket("/chat")
workflow StreamingChat(message: string, session_id: string) -> stream<string> {
    return Assistant.stream(message, session: session_id)
}

fn main() {
    server = http.Server([StreamingChat])
    server.listen(8080)
}
\end{lstlisting}

The \code{stream<string>} return type tells \haira{} to stream each token over the WebSocket as it arrives. No buffering --- the client sees tokens in real time.

\section{Dual-Mode Chatbot}

Expose the same agent over both REST (for simple integrations) and WebSocket (for real-time UIs):

\begin{lstlisting}
agent Assistant {
    provider: anthropic
    system: "You are a helpful assistant."
    memory: conversation(max_turns: 50)
}

@webhook("/chat")
workflow ChatREST(message: string, session_id: string) -> { reply: string } {
    reply, err = Assistant.ask(message, session: session_id)
    return { reply: reply or "Error occurred." }
}

@websocket("/chat/stream")
workflow ChatStream(message: string, session_id: string) -> stream<string> {
    return Assistant.stream(message, session: session_id)
}

fn main() {
    server = http.Server([ChatREST, ChatStream])
    server.listen(8080)
}
\end{lstlisting}

Both endpoints share the same agent and session memory.

\section{RAG Chatbot}

A chatbot that searches a knowledge base before answering:

\begin{lstlisting}
import "tools/postgres"
import "tools/http"

tool search_kb(query: string) -> [KBArticle] {
    """Search the knowledge base for relevant articles"""

    results, err = pg.query(
        "SELECT * FROM articles WHERE content @@ to_tsquery($1) LIMIT 5",
        [query]
    )
    if err != nil { return [], err }
    return results
}

tool lookup_order(order_id: string) -> Order {
    """Look up a customer order by ID"""

    order, err = pg.query_one(
        "SELECT * FROM orders WHERE id = $1",
        [order_id]
    )
    if err != nil { return Order{}, err }
    return order
}

tool create_ticket(subject: string, body: string, priority: string = "normal") -> Ticket {
    """Create a support ticket"""

    resp, err = http.post("https://helpdesk.api/tickets", {
        subject: subject, body: body, priority: priority
    })
    if err != nil { return Ticket{}, err }
    return json.decode(resp.body, Ticket)
}

agent SupportBot {
    provider: anthropic
    system: """
        You are a customer support agent for Acme Corp.
        Use the knowledge base to answer questions.
        Look up orders when customers ask about purchases.
        Create tickets for issues you cannot resolve directly.
        Always be polite and helpful.
    """
    tools: [search_kb, lookup_order, create_ticket]
    memory: conversation(max_turns: 100)
    temperature: 0.3
}

@websocket("/support")
workflow CustomerSupport(message: string, session_id: string) -> stream<string> {
    return SupportBot.stream(message, session: session_id)
}
\end{lstlisting}

The agent automatically decides when to search the knowledge base, look up an order, or create a ticket based on the user's message.

\section{Multi-Agent Routing}

Route user messages to specialized agents based on intent:

\begin{lstlisting}
agent Router {
    provider: anthropic
    system: """
        Classify the user message into one category: billing, technical, general.
        Reply with just the category name, nothing else.
    """
    temperature: 0.0
}

agent BillingAgent {
    provider: anthropic
    system: "You handle billing questions. You can look up invoices and process refunds."
    tools: [lookup_invoice, process_refund]
    memory: conversation(max_turns: 30)
}

agent TechAgent {
    provider: anthropic
    system: "You handle technical support. You can search docs and check system status."
    tools: [search_docs, check_status]
    memory: conversation(max_turns: 30)
}

agent GeneralAgent {
    provider: anthropic
    system: "You handle general inquiries."
    memory: conversation(max_turns: 30)
}

@websocket("/chat")
workflow SmartChat(message: string, session_id: string) -> stream<string> {
    category, _ = Router.ask(message)

    match category {
        "billing"   => return BillingAgent.stream(message, session: session_id)
        "technical" => return TechAgent.stream(message, session: session_id)
        _           => return GeneralAgent.stream(message, session: session_id)
    }
}
\end{lstlisting}

\section{Agent Handoffs}

For complex support bots, use handoffs instead of explicit routing. The front desk agent decides when to delegate, and the runtime handles the transfer automatically:

\begin{lstlisting}
agent FrontDesk {
    provider: anthropic
    system: """
        You are a front desk agent. Greet users and help with general questions.
        If they have a billing question, hand off to BillingAgent.
        If they have a technical issue, hand off to TechAgent.
    """
    handoffs: [BillingAgent, TechAgent]
    memory: conversation(max_turns: 10)
}

agent BillingAgent {
    provider: anthropic
    system: "You handle billing questions. You can look up invoices and process refunds."
    tools: [lookup_invoice, process_refund]
    memory: conversation(max_turns: 30)
}

agent TechAgent {
    provider: anthropic
    system: "You handle technical support. You can search docs and check system status."
    tools: [search_docs, check_status]
    memory: conversation(max_turns: 30)
}

@webhook("/chat")
workflow SmartSupport(message: string, session_id: string) -> { reply: string } {
    // FrontDesk automatically hands off when appropriate
    reply, err = FrontDesk.ask(message, session: session_id)
    if err != nil { return { reply: "Sorry, something went wrong." } }
    return { reply: reply }
}

fn main() {
    server = http.Server([SmartSupport])
    io.println("Support server on :8080")
    server.listen(8080)
}
\end{lstlisting}

Compare this to the explicit routing in the previous section. With handoffs, the workflow code stays simple --- the agent decides when to delegate, not the workflow. This is ideal for conversational bots where the routing logic is nuanced and best handled by the LLM itself.

\section{Chatbot with Structured Actions}

Sometimes you want the chatbot to return structured data alongside text:

\begin{lstlisting}
struct ChatResponse {
    reply: string
    actions: [Action]
}

struct Action {
    type: string      // "navigate", "show_modal", "refresh"
    target: string
}

agent UIAssistant {
    provider: anthropic
    system: """
        You are a UI assistant. When answering, also suggest UI actions.
        For example, if the user asks about their order, suggest navigating
        to the orders page.
    """
    tools: [search_kb, lookup_order]
    memory: conversation(max_turns: 50)
}

@webhook("/assistant")
workflow UIChat(message: string, session_id: string) -> ChatResponse {
    response: ChatResponse, err = UIAssistant.run(
        { message: message },
        session: session_id
    )
    if err != nil {
        return ChatResponse{ reply: "Sorry, an error occurred.", actions: [] }
    }
    return response
}
\end{lstlisting}

\section{Best Practices}

\begin{enumerate}
    \item \textbf{Always use sessions for conversation} --- without a session, memory agents start fresh every call
    \item \textbf{Set appropriate max\_turns} --- too high wastes tokens, too low loses context
    \item \textbf{Use streaming for interactive UIs} --- REST is fine for backend-to-backend, but users expect real-time responses
    \item \textbf{Give tools clear descriptions} --- the agent decides when to use tools based on descriptions
    \item \textbf{Use a Router agent for complex bots} --- instead of one agent doing everything, route to specialists
    \item \textbf{Handle errors gracefully} --- always provide a fallback response when agents fail
\end{enumerate}
